\section{The one-site comparison and susceptibility calibration}\label{sec:cavity}

We prove the local scalar closure used below. The removal of one site
belongs to the cavity framework of \cite{AizenmanSimsStarr2003}; here
its quantitative remainder is controlled by a finite family of test
replicas and a susceptibility calibration. All estimates in this
section concern a minimizer of the changed paths \eqref{eq:6}
in a direct test. We assume outer regularity above a fixed multiple
of \(x\), and
\[
             x\gtrsim a,\qquad e(x)x^2\gtrsim nE^2,
                                      \qquad e(x)\gtrsim x^2.
\]
These are the hypotheses for the direct-test consequences of
\eqref{eq:18}. In particular,
\[
 W=\max\{x,\sqrt{e(x)}\},\qquad x\le W\lesssim d,\qquad x\longrightarrow0.
\]
We also assume the endpoint hypothesis \eqref{eq:4}
throughout the regular outer region. More explicitly,
let \(C_{\rm reg}\) be a fixed outer-regularity threshold.
Along every sequence of masses \(u\) in the open working
interval with \(u\ge C_{\rm reg}x\), we assume
\[
                  b_0(u)=T+o(u^2),\qquad h_0(u)=o(u^3).
\]
The little-oh estimates are uniform on fixed comparable
mass bands, in the sequential sense specified after
\eqref{eq:4}. In particular they hold on the regular
calibration bands of scale \(W\) used below. No instance
of \eqref{eq:4} is imposed at a target or omitted lower
cut below this regular region.

Constants may depend on the fixed temperature, fixed outer
comparability constants, and the fixed number of test factors, but not
on the number of steps in a step approximation. An estimate for an
integrated replica tree means an estimate after integration against
the total variation of its deterministic attachment measure. It does
not assert pointwise concentration at every high split.

Write \(b,h\) for the changed paths and put \(K=h+bS\).
The path \(K\) is nonnegative and nondecreasing. Denote by
\(\Gamma(K_i)\) the second conditional-magnetization moment at cut
\(i\) of the scalar Ising recursion with the entire covariance path
\(K\). This notation retains dependence on the entire path, not
merely on the number \(K_i\).

\begin{proposition}[Uniform one-site closure]\label{prop:cavity-closure}
Fix \(0<c<C<\infty\). At a target cut \(i\) satisfying
\(r_i\le Cx\) and \(cx\le u_i\le Cx\), along every sequence with
\(D_i=o(x)\),
\begin{equation}\label{eq:22}\tag{22}
                         S_i=\Gamma(K_i)+o(x^3).
\end{equation}
The bound is \(O(x^3)\), uniformly over these targets, if \(D_i=O(x)\).
More precisely, for a fixed direct-test sequence there are
\(\eta_n\to0\), \(C_0<\infty\), and \(\vartheta>0\), independent of
the target in this band, such that
\[
 \frac{|S_i-\Gamma(K_i)|}{x^3}
       \le C_0\{\eta_n+(D_i/x)^\vartheta\}.
\]
When \(W=O(x)\), the \(O(x^3)\) assertion also holds at every omitted
lower cut, without a lower bound on its mass and without
\eqref{eq:4} at that cut. On the bands in \eqref{eq:19},
the error in \eqref{eq:22}, divided by \(x^3\), tends to zero in
normalized \(L^1\), in both mass and radius coordinates.
\end{proposition}

The two scales in this statement play different roles. At the target
scale \(x\), the comparison gives integrated control of overlap errors,
but this alone does not control the susceptibility correction left by
removing one site. When \(W=o(1)\), a regular split of scale \(W\)
has a centered overlap that is polynomially small in every fixed
moment. We test the one-site identity at that split and use each new
test factor to gain another such saving. A fixed number of tests then
controls the susceptibility correction. A second mean identity at
scale \(x\) transfers this control to the target cut. When \(W\)
stays bounded below, the integrated overlap bound controls the
correction directly.

\subsection{The calibration scale and its test factors}

We record the preceding estimates that will be used. At every fixed
moment order the raw bounds \eqref{eq:12} hold. For \(W=o(1)\), choose
a good cut \(*\) of scale \(W\), and a continuity cut of mass
\(U\asymp W\) at a sufficiently large fixed multiple of that scale.
Both are regular, with a regular innovation interval between them.
The good cut can be chosen above any prescribed fixed multiple of
\(x\). Equations \eqref{eq:17}--\eqref{eq:19} give
\[
 S_*\asymp W,\qquad
 D_*\le n^{-\kappa}\frac{x^2}{W},\qquad
 \|R_*-S_*\|_p\le C_p n^{-c_2}W,\qquad
 \int_0^{C_1W}D_u\,du=o(x^2).
\]
The exponent \(c_2>0\) is independent of each fixed \(p\);
\(\kappa>0\) can be decreased after finitely many fixed-order
H\"older estimates. At a noncoincident high error split, integrated
with a bounded split density above \(U\), \eqref{eq:18} gives a
polynomially saved \(x^2/W\) bound in \(L^2\). The same bound holds
at \(*\). Deleted-coordinate overlaps satisfy these statements with
an additional \(1/n\).

A block is a set of leaves sharing through the cut of mass \(U\).
A boundary edge joins distinct such blocks. If its low split is at
most a fixed multiple of \(x\), or is \(*\), changing an endpoint
by continuation resampling after a regular cut of scale \(v\ge U\)
costs, by \eqref{eq:16},
\[
 C_p n^{-\kappa}z_v,\qquad z_v=\frac{x^3}{W^{3/2}\sqrt v}.
\]
For a projected edge anchored by \(X_*\) the bound is
\[
                         C_p n^{-\kappa}\frac{x^3}{Wv}.
\]
Indeed, \eqref{eq:13}, on an innovation interval before the high
cut and after the low fork, bounds this projection by
\(C_p(\log n)^C\sqrt W/\sqrt{n e(v)v^3}\).
Use \(e(v)v\gtrsim W^3\) and \(n^{-1/2}\ll x^3\).
For an unprojected boundary edge use the synchronized two-branch
decomposition in \eqref{eq:16}. These estimates hold for the
rest-only system below: its external-field clock is unchanged,
and its raw moments transfer by Lemma~\ref{lem:cavity-row-change}.

Throughout, \(\nu\) denotes normalized planted expectation for the
actual changed path on the prescribed replica tree. For a noncoincident
pair of leaves \(c,d\), let \(j(c,d)\) be its prescribed shared
revelation cut, including the specified side at a jump, and write
\[
 \Delta_{cd}=R_{cd}-S_{j(c,d)},\qquad
 b_{cd}=b_{j(c,d)},\qquad K_{cd}=K_{j(c,d)}.
\]
All centers \(S_j\) are those of the actual changed path and remain
fixed during row interpolation. Coincident-label contributions retain
the terminal and diagonal conventions of \eqref{eq:2}.

\begin{definition}[Admissible tests]\label{def:cavity-tests}
A test of degree at most \(N\) is a finite linear combination of
products of centered boundary overlaps and their conditional
projections, with these specifications. Increase \(N\), if necessary,
so that it also bounds the number of leaves after the fresh-replica
representation of the projections. All low splits are \(*\),
except for at most one initial edge at a cut \(i\).
The distinguished target \(a\) is unique in its block among the
observed test leaves. Its dependence is a power of
\[
                         Y_a=X_*\cdot y_a-S_*
\]
times at most one incoming boundary edge from an outside leaf.
Other factors do not involve \(a\). The second target \(b\), when
an identity is tested, is fresh at \(*\) and absent from the test.

An old block may contain several high relatives. Its internal
relationships occur only in deterministic scalar coefficients and
the prescribed tree topology; no overlap factor is internal to a
cut block. High relationships between old leaves are integrated
with finite-label attachment measures. Their bounded scalar
coefficients, and constants obtained by expectations of earlier
tests, are fixed during a subsequent row interpolation.
For a monomial, its natural scale is the product of \(W\) for each
star factor and \(x\) for the possible initial factor, including
factors placed in a separate expectation. For a linear combination,
\(p_0\) is the sum of these scales multiplied by the absolute bounds
on their scalar coefficients. The degree and the number of leaves
are fixed independently of \(n\); all tests generated below have
these bounds controlled by the initial degree and recursion depth.
\end{definition}

Represent projection powers by conditionally independent fresh
teeth at their prescribed low split. Tests therefore become
spin-only tests before interpolation. Every assertion below is for
a fixed \(N\), with constants allowed to depend on \(N\).

\begin{lemma}[Bounded row changes]\label{lem:cavity-row-change}
Along a segment of admissible changes of one Gaussian row with
uniformly bounded covariance kernels, expectations of nonnegative
spin-only tests on at most \(N\) leaves change by at most a factor
\(C_N\), independently of the step partition. The raw moment
bounds and the actual-path centered-overlap bounds above therefore
hold throughout such a segment with this factor.
\end{lemma}

\begin{proof}
Apply \eqref{eq:2}. The row kernel is bounded and the total
variation of its two added-label measure is bounded in terms of
\(N\). Marginalizing the new leaves leaves the original test,
so \(|\partial_t\nu_t[A]|\le C_N\nu_t[A]\). Gr\"onwall in both
directions proves the assertion. Centering constants are those
of the actual path and are fixed. When the row is scalar, its
recursion factors from the rest-spin recursion. The latter
retains all original external-field increments on the remaining
coordinates. The proofs of \eqref{eq:13} and \eqref{eq:16} use
these increments and raw moments, so apply to this rest-only law.
No identification of interpolated and actual conditional means
is needed.
\end{proof}

\subsection{Attachment cancellation with observed high subtrees}

The signed differentiation rule is summed before taking absolute
values. We specify its cancellation inside an observed block.

\begin{lemma}[Truncation inside an observed block]\label{lem:cavity-truncation}
Fix an old tree with at most \(N\) leaves in a block of mass \(U\).
Deterministic scalar coefficients may depend on all its old
split values and earlier test expectations, but not on a new
leaf \(c\)'s deeper placement. Replace \(c\), after an occupied
prefix of mass \(v\), by an independent continuation from that
prefix. The signed sum of deeper placements having this same
truncated value is \(v\) times that value. The telescoping
difference between truncations at adjacent comparable scales
\(v,2v\) has total variation at most \(C_Nv\). The final layer
may include coincidences.
\end{lemma}

\begin{proof}
The exact collapse is Lemma~\ref{lem:found-attachment}, and the
layer estimate is Lemma~\ref{lem:found-cut-truncation}. Their
coefficient restriction holds because the old topology and earlier
test expectations are fixed before \(c\) is introduced.
For two old leaves splitting at \(v_B\) inside a block of
mass \(U\), for example, the total is
\[
                  2-2(1-v_B)-v_B-(v_B-U)=U.
\]
Here \(-v_B\) is the old-fork atom. For coincident leaves
the corresponding identity is \(1-(1-U)=U\).
These are signed identities; replacing the coefficients by their
absolute values before this grouping would lose the factor \(U\).
\end{proof}

\begin{lemma}[Boundary covariance estimate]\label{lem:cavity-boundary}
Let \(F_c\) be a product of two factors of the form
\(h_j+b_jR_{cd}\), where \(d\) is exterior to the block of
\(c\) and their low split satisfies \(j=j(c,d)\in\{i,*\}\);
one exterior endpoint
may be conditionally projected. Let \(P\) be admissible.
If no overlap edge of \(P\) is
internal to the old block in which \(c\) is allocated,
the placement sum after subtracting the fresh-at-cut
baseline is \(o(Wx^2p_0)\), uniformly at fixed degree.
\end{lemma}

\begin{proof}
Apply Lemma~\ref{lem:found-coupling} to each grouped layer of
Lemma~\ref{lem:cavity-truncation}. Its conditioning consists of the
complete prefix of the affected occupied subblock and all branches
that have already departed. The new leaf supplies \(F_c\); the
affected old subtree supplies \(P\). The other endpoints of the
two factors in \(F_c\), and of every varying edge in \(P\), lie
outside that subtree. Before an old fork we retain both old children;
after it the other child belongs to the exterior. The normalized
transition kernels therefore give exactly the two conditional
marginals required by the coupling lemma, including for coincidence.
Its centered product estimate applies after averaging the prefix.

Lemma~\ref{lem:found-boundary-tests}, the resampling estimates above,
and the raw moments now give
\[
 \|F_c-F_c'\|_p\le C_{N,p}n^{-\kappa}Wz_v,\qquad
 \|P-P'\|_p\le C_{N,p}n^{-\kappa}(z_v/x)p_0.
\]
The second bound uses the smallest factor scale \(x\);
star factors give the stronger denominator \(W\).
Lemma~\ref{lem:cavity-truncation} bounds each layer by a
polynomially saved multiple of
\[
             v(Wz_v)(p_0z_v/x)=p_0x^5/W^2\le p_0Wx^2.
\]
There are \(O(\log n)\) layers on these polynomial scales.
Their number is absorbed by the saving. This argument
is uniform in the old topology and its split values,
so bounded-variation integration preserves the estimate.
\end{proof}

\subsection{The row expansion}\label{subsec:cavity-row-expansion}

Distinguish one coordinate \(k_0\) with spin \(s\), so
\(s_c=\sigma_{k_0}^c\), and normalize deleted overlaps by the full size
\(n\):
\[
 R^-_{cd}=\frac1n\sum_{k\ne k_0}\sigma_k^c\sigma_k^d
         =R_{cd}-\frac{s_cs_d}{n}.
\]
Its row covariance between distinct leaves \(c,d\) at shared cut
\(j=j(c,d)\) is
\[
 s_cs_dQ_{cd},\quad Q_{cd}=h_j+b_jR^-_{cd},\quad
 Q_{cd}-K_j=b_j\Delta^-_{cd},\quad
 \Delta^-_{cd}=R^-_{cd}-S_j.
\]
All centers remain those of the actual path. Replacing
a deleted overlap by the full overlap in a degree-\(N\)
test costs \(C_Np_0/(nx)\). Permutation symmetry then
replaces a distinguished-site factor by its overlap
average at the actual model with the same error.
These are harmless since
\[
                  (nx)^{-1}=o(x^3),\qquad (nx)^{-1}=o(Wx^2),
\]
using \(nx^4\to\infty\) and \(W\ge x\).
Common diagonal row noise cancels by \eqref{eq:2}.

First suppose \(W=o(1)\). Below \(U\) multiply the row
covariance by \(l\in[0,1]\), moving missing revelation
to \(U\). Before this operation mix the scalar row
\(K\) with the actual row \(Q\), in fraction \(t\).
Write \(\nu_{l,t}\) for the resulting admissible law.
For each \(l\in[0,1]\), let \(\nu_l^s\) be the scalar-site marginal
of \(\nu_{l,0}\). Its cumulative clock is \(lK_u\) below \(U\) and
\(K_u\) above \(U\), with the missing revelation supplied at mass
\(U\) and the specified revelation sides retained. Write
\(\nu^s=\nu_0^s\); the law \(\nu_1^s\) uses the full scalar path
\(K\). Unadorned \(\nu\) is the actual changed-path expectation,
recovered as \(\nu_{1,1}\). On every fixed replica tree, the scalar-site
and remaining-coordinate recursions factor at \(t=0\).
At \(l=0\) the site has independent flip symmetries
in its cut blocks. At \(l=t=0\) it is an independent
scalar recursion.
For noncoincident scalar leaves sharing to \(v\in(U,1)\), write
\(\gamma(v)=\nu^s[s_as_c]\), retaining the specified side at a jump;
for coincidence, \(\gamma_{aa}=1\). All free-label sums use the signed
attachment measures of \eqref{eq:2} and
Lemma~\ref{lem:found-attachment}, with the displayed restrictions on
coincidences. For the block \(A\) of \(a\), put
\[
             \gamma_{ac}=\nu^s[s_as_c],\qquad
             \chi=\sum_{c\in A}\gamma_{ac}.
\]
The scalar row sum is independent of other observed
labels. It is the scalar continuation Hessian at zero incoming field
before the relocated increment at \(U\), with that increment included
in the continuation:
\[
       \chi=1-\int_U^1\gamma(v)\,dv,\qquad 0<c_0\le\chi\le C_0.
\]
Its lower bound follows from expected terminal thermal
variance, since total scalar clock and planted drift
are bounded. For distinct \(a,c\) sharing to high mass
\(v\), \(\gamma_{ac}\le Cv\): the clock is \(O(v)\),
the magnetization martingale starts at zero, and its
spatial derivatives are bounded.

Define the linear insertion operators
\begin{equation}\label{eq:23}\tag{23}
 \begin{split}
 G_{ac}&=\frac12\sum_{\substack{e,f\in A\\e\ne f}}
 \{\nu^s[s_as_cs_es_f]-\gamma_{ac}\gamma_{ef}\}
                                      b_{ef}\Delta_{ef},\\
 M_a&=\sum_{c\in A}G_{ac}.
 \end{split}
\end{equation}
These are interpreted inside a test pairing and planted
expectation, as signed insertions. In particular
\(\nu[M_a]=0\), because each noncoincident actual-path
centered overlap has zero marginal expectation.

The test family permits two operations: adding a projected star
factor at the current target, and choosing a fresh target joined by
a star edge to a high relative of the former target. The next lemma
verifies that both operations preserve the boundary structure needed
in Lemma~\ref{lem:cavity-boundary}.

\begin{lemma}[Closure of the test family]\label{lem:cavity-family}
Centering an admissible test, multiplying by a projected
star factor at its target, or passing to a fresh target
and multiplying by one centered star edge to a summed
high relative in the former target block preserves
Definition~\ref{def:cavity-tests} at a larger finite degree.
New coefficients and attachment measures are bounded in
terms of the resulting degree.
\end{lemma}

\begin{proof}
Regard cut blocks as vertices and overlap factors as
edges. Each edge joins different vertices. A projection
adds fresh teeth at \(*\), hence boundary edges.
A switch adds an edge from a new vertex to a high
relative in the former target vertex. It adds no overlap
edge between that relative and the former target:
their relation enters only through a deterministic
scalar coefficient and the old topology. The new target
is unique in its block and absent from previous factors.
Centering replaces monomials by earlier expectations,
independent of subsequent leaf placements.

Two switches from \(\Delta_{ao}\), for example, give
\[
 (\Delta_{ao}\Delta_{bc}-\nu[\Delta_{ao}\Delta_{bc}])
                                                    \Delta_{de}
\]
with old scalar coefficient \(\gamma_{ac}\gamma_{be}\).
In the block containing \(b,e\) the varying edges are
\(bc,de\); there is no \(be\) overlap factor. Before
their old fork resample \(b,e\) jointly, retaining
the fork. Test variation is \(O_p(xWz_v)\), at most
\(O_p(p_0z_v/x)\). The coefficients are fixed by the
old topology, so Lemma~\ref{lem:cavity-truncation}
applies. Induction gives the same graph property
for every fixed number of switches. Finite-label
total-variation and split-density bounds from
\eqref{eq:2} depend on this degree, not the mesh
or the least mass.
\end{proof}

\begin{lemma}[Finite-degree row identity]\label{lem:cavity-row-identity}
Let \(a,b\) split at \(*\), with \(b\) fresh, and let
\(P\) be admissible of degree \(N\). Let \(A\) and \(B^{\prime}\)
be the blocks at the current cut \(U\) containing \(a\) and \(b\),
respectively. Then
\begin{equation}\label{eq:24}\tag{24}
 \begin{split}
 \nu[R_{ab}P]={}&\Gamma(K_*)\nu[P]
  +b_*\sum_{c\in A,d\in B'}\gamma_{ac}\gamma_{bd}
                                             \nu[\Delta_{cd}P]\\
 &+\sum_{c\in A,d\in B'}
   \nu[P(G_{ac}\gamma_{bd}+\gamma_{ac}G_{bd})
                                     (K_*+b_*\Delta_{cd})]
   +\mathcal R_N(P).
 \end{split}
\end{equation}
At fixed degree, uniformly over initial targets \(i\) in the band
of Proposition~\ref{prop:cavity-closure}, the remainder satisfies
\[
 |\mathcal R_N(P)|\le
 C_NWx^2p_0\{\eta_{N,n}+(D_i/x)^{\vartheta_N}\},\qquad
       \eta_{N,n}\to0,\quad\vartheta_N>0,
\]
where the \(D_i\) term is absent without an initial factor.
Consequently \(\mathcal R_N(P)=o(Wx^2p_0)\) for star-only tests
and along every moving initial target with \(D_i=o(x)\).
\end{lemma}

\begin{proof}
We first control the third remainder in the low revelation parameter.
For the first and second low coefficients we then freeze the actual
low kernel \(Q\) while expanding only the high law; this is a separate
expansion from the mixed derivative with its varying kernel \(L_t\).
It isolates the retained linear correction in \eqref{eq:24} and a
second-order term passing through a third block. The latter is treated
separately for an unobserved and an observed block.
First delete the distinguished site from every test.
Let \(\mathcal I_\mu[J;V_1,\ldots,V_k]\) denote
\(k\) signed covariance insertions, including
\(2^{-k}\) from \eqref{eq:2}. Kernels include site spins.
Write
\[
 \delta_{cd}=b_j\Delta^-_{cd},\quad
 L_t=(s_cs_d(K_j+t\delta_{cd}))_{\rm low},\quad
 E_l=(s_cs_d\delta_{cd})_{\rm high}
                      +l(s_cs_d\delta_{cd})_{\rm low}.
\]
Here low means distinct cut blocks. Set
\[
 F(l,t)=\nu_{l,t}[s_as_bP]-g(l)\nu_{l,t}[P],\qquad
 g(l)=\nu_l^s[s_as_b].
\]
In particular, \(g(1)=\Gamma(K_*)\) for the full scalar path.
Then \(F(0,t)=F(l,0)=0\). Since the low covariance is
\(l(K+t\delta)\), the exact differentiation rule, for an integer
\(k\ge0\), is
\[
 \partial_t\partial_l^k\nu_{l,t}[J]
 =\mathcal I_{\nu_{l,t}}[J;E_l,L_t^{\otimes k}]
  +k\mathcal I_{\nu_{l,t}}
                    [J;\delta_{\rm low},L_t^{\otimes(k-1)}],
\]
the second term being absent for \(k=0\).

\emph{Third low remainder.}
The third low Taylor remainder in \(F(1,1)\) is
\[
 \frac12\int_0^1(1-l)^2\int_0^1 F_{lllt}(l,t)\,dt\,dl .
\]
Its product term uses the full Leibniz rule
\[
 \partial_t\partial_l^3(g\nu[P])
   =\sum_{r=0}^3\binom3r g^{(r)}
                  \partial_t\partial_l^{3-r}\nu[P].
\]
Thus each measure-derivative term has one error and
three low factors, and each mixed-covariance derivative
has one low error and two low factors. A scalar
\(g^{(r)}=O_N(W^r)\) counts as \(r\) low factors.
The mixed terms have multiplicity
\(\binom3r(3-r)\), in addition to insertion normalizations.

Allocate an error pair before allocating the other
new labels. A noncoincident high error has a bounded
split density, including old high splits after their
old integration, and costs a polynomially saved
\(x^2/W\). So does an error at \(*\).
A new low error has bounded split density and costs
\(\int_0^{C_1W}D=o(x^2)\). This order of allocation
prevents a newly created fork from becoming a fixed
atom. The only other old low atom is \(i\).

At \(i\), an ordinary measure-side edge crossing a
regular cut of scale \(Cx\) above \(i\) and below
\(*\) costs \(O_p(x)\). If there is none, turn off
site revelation below this smaller cut. The
\(i\)-error is odd in its two blocks; the target
factor, if present, is even and other measure-side
insertions remain inside blocks. The expectation
vanishes. Restoring revelation by one derivative
supplies an \(O_p(x)\) edge. Scalar \(g^{(r)}\)
factors cannot repair this parity. In particular
the term \(3g''\mathcal I[P;\delta_{\rm low}]\),
which has no ordinary low edge in its measure,
costs at most
\[
 C_NW^2x(D_i+1/n)p_0\le C_NWx(D_i+1/n)p_0.
\]
Every old-\(i\) term therefore has the latter bound,
up to fixed-moment interpolation of \(D_i/x\).
If \(W=O(x)\), its two original low factors suffice.
Diagonal errors vanish. This treats every Leibniz
and mixed-covariance term of the third remainder.

\emph{Expansion with the actual low kernel fixed.}
For the first and second low coefficients, freeze
the actual low kernel \(Q\) while changing only the
high law. Put \(\mu_t=\nu_{0,t}\), \(c_j=g^{(j)}(0)\),
and
\[
 \begin{split}
 A(t)&=\mathcal I_{\mu_t}[s_as_bP;Q_{\rm low}]
                                              -c_1\mu_t[P],\\
 B(t)&=\mathcal I_{\mu_t}[s_as_bP;Q_{\rm low},Q_{\rm low}]
       -2c_1\mathcal I_{\mu_t}[P;Q_{\rm low}]-c_2\mu_t[P].
 \end{split}
\]
In particular \(A(t)\ne F_l(0,t)\) in general.
The exact decomposition is
\[
 \begin{split}
 F(1,1)={}&A(0)+A'(0)+\tfrac12B(0)
       +\int_0^1(1-t)A''(t)\,dt
       +\tfrac12\int_0^1B'(t)\,dt\\
 &+\frac12\int_0^1(1-l)^2\int_0^1F_{lllt}(l,t)\,dt\,dl .
 \end{split}
\]
The \(A''\) terms have two high errors and one low
factor; the \(B'\) terms have one high error and
two low factors. Their bounds are respectively
\[
 C_Nn^{-\kappa}W(x^2/W)^2p_0
       \le C_Nn^{-\kappa}Wx^2p_0,\qquad
 C_Nn^{-\kappa}W^2(x^2/W)p_0.
\]
For two saved \(L^2\) errors use H\"older exponents
just above 2, interpolating with bounded moments;
the polynomial saving pays the fixed loss.
A lone low error uses Cauchy--Schwarz and raw
higher moments, requiring no rate for its little-oh.

\emph{The retained first-order terms.}
Parity restricts the low edge in \(A(0)\) to
\(c\in A,d\in B'\), and \(c_1=\chi^2K_*\).
Define
\[
 C(t)=b_*\sum_{c\in A,d\in B'}\gamma_{ac}\gamma_{bd}
                                      \mu_t[\Delta^-_{cd}P].
\]
Then \(A(0)=C(0)\). Subtracting \(C'(0)\) from
\(A'(0)\) leaves the high insertion coefficient
\[
 \nu^s[s_as_cs_bs_ds_es_f]
                         -\gamma_{ac}\gamma_{bd}\gamma_{ef}
\]
times \(Q_{cd}P\). Scalar independence cancels it
unless the high edge lies in \(A\) or \(B'\).
There it is
\((G_{ac}\gamma_{bd}+\gamma_{ac}G_{bd})Q_{cd}\).
Replacing \(C(0)+C'(0)\) by \(C(1)\) adds two
high errors and a low factor. Restoring the
rest-only measure in the connected insertion
adds one further high error. Both have the
\(A''\) bound above.

Restoring full low revelation in \(C(1)\), the
first derivative at zero vanishes by block parity;
the second costs \(C_ND_*W^2p_0=o(Wx^2p_0)\).
For the connected high insertion an absolute first
derivative costs one high error and two low
factors, covered by the \(B'\) bound. These are
the actual-measure terms in \eqref{eq:24}.

\emph{The second low coefficient.}
It remains to bound \(B(0)/2\). Its single-low
insertion vanishes by parity. Two low edges paired
with \(s_as_b\) must form a path \(A-D-B'\)
through a third block \(D\). Two edge orders and
two orientations per edge, times \(1/4\) from
the insertions, give exactly the coefficient 2.
The scalar \(c_2\) cancels this same sum with
\(QQ\) replaced by \(KK\). The Taylor factor
\(1/2\) leaves exactly
\[
 \sum_{D\ne A,B'}\ \sum_{\substack{f\in A,\ c,e\in D\\d\in B'}}
 \gamma_{af}\gamma_{ce}\gamma_{db}
          \mu_0[P\{Q_{fc}Q_{ed}-K_{fc}K_{ed}\}].
\]
This includes \(c=e\), \(f=a\), and \(d=b\);
none is a diagonal low edge.

\emph{Unobserved third block.}
If \(D\) is unobserved, its first entry leaves the
old tree below \(U\), with bounded departure
density or old-fork coefficient \(O(W)\).
Use \(QQ-KK=\delta Q+K\delta\).
A varying low error costs \(o(x^2)\) after
integration, while the other factor costs \(W\).
At \(*\) the costs are
\(o_{\rm poly}(x^2/W)\), \(W\), and the entry
coefficient \(O(W)\). If an edge splits at \(i\),
both do; the product difference is
\(O(x^2(D_i/x)^{\vartheta_N})\), before entry.
A departure along a terminal trunk between
\(*\) and \(U\) has one varying split and the
other fixed at \(*\), covered by the same
bounds. These exhaust a new third block's
locations relative to the target pair.

\emph{Observed third block.}
If \(D\) is observed, its two exterior splits
are both \(*\) or both \(i\). Average the fresh
endpoint \(d\) to its projection at \(*\),
and sum its scalar row to \(\chi\). This remains
valid in the latter case. Evaluate the factor
depending on \(e\) at \(c\). For distinct \(c,e\)
at high split \(v\), its error is a saved
\(x^3/(Wv)\). The factor \(\gamma_{ce}\le Cv\)
pays this denominator, including at old high
atoms; the other low factor costs at most \(W\).
The result is \(o(x^3p_0)\le o(Wx^2p_0)\).
At coincidence the error is zero. The resulting
rest observable is independent of \(e\)'s deeper
placement, so its scalar row sums to \(\chi\).

Only unweighted allocation of \(c\) in \(D\)
remains. Its fresh-at-cut baseline gains \(U\).
The baseline product difference is \(O(WD_*)=
o(x^2)\) at \(*\), or
\(O(x^2(D_i/x)^{\vartheta_N})\) at \(i\).
Restoring deeper attachment is precisely
Lemma~\ref{lem:cavity-boundary}. By
Lemma~\ref{lem:cavity-family} its hypotheses
hold also for several old high relatives
and earlier centering constants.

For completeness the exhaustive remainder
ledger, with \(p_0\) suppressed, is
\begin{center}
\small
\begin{tabular}{p{0.29\linewidth}p{0.29\linewidth}p{0.30\linewidth}}
\toprule
Class & Factors & Bound\\
\midrule
Third low: measure derivative &
one error, three low factors &
\(o(Wx^2)+O(WxD_i)\)\\
Third low: mixed covariance &
one low error, two low factors &
\(o(Wx^2)+O(WxD_i)\)\\
First high remainder and high restoration &
two high errors, one low factor &
\(o_{\rm poly}(x^4/W)\)\\
Second high remainder and connected low restoration &
one high error, two low factors &
\(o_{\rm poly}(Wx^2)\)\\
Centered first low restoration &
star error, two low factors &
\(O(D_*W^2)\)\\
Scalar-high second, new third block &
centered product; departure or entry &
\(o(Wx^2)+O(Wx^2(D_i/x)^{\vartheta_N})\)\\
Scalar-high second, old third block &
row reduction, baseline, covariance &
\(o(Wx^2)+O(Wx^2(D_i/x)^{\vartheta_N})\)\\
\bottomrule
\end{tabular}
\end{center}
All derivatives occur in the exact decompositions
above. Allocating an error pair first gives only
diagonal errors, old low atoms, integrated old
high splits, or new splits of bounded density.
Remaining allocations have bounded total
variation at fixed degree. Interpolate
\(\Delta_i/x\) between its \(L^2\) norm and
bounded higher moments if a higher H\"older
exponent is required. The finitely many
resulting positive exponents have a positive
minimum \(\vartheta_N\), proving the claimed
uniform modulus. Finally restore deleted
overlaps and use permutation symmetry.
\end{proof}

\subsection{Projected replacements and the longitudinal insertion}

Write \(\mathcal L(P,a)=\nu[PM_a]\). The replacements needed
to close \eqref{eq:24} are justified next.

\begin{lemma}[Projected replacements]\label{lem:cavity-replacements}
After averaging the fresh partner block,
\[
 \sum_{c\in A}\gamma_{ac}\nu[Y_cP]
              =\chi\nu[Y_aP]+o(Wx^2p_0).
\]
In the centered part multiplying \(G_{ac}\), \(Y_c\) can
likewise be replaced by \(Y_a\), with error \(o(Wx^2p_0)\).
In the part multiplying \(G_{bd}\), interchanging \(b,d\)
and summing that block gives exactly
\(\sum_c\gamma_{ac}\mathcal L(P\Delta_{bc},b)\).
The uniform target modulus of
Lemma~\ref{lem:cavity-row-identity} applies.
\end{lemma}

\begin{proof}
The partner operation is exact: its entire block is fresh,
\(P\) does not use its leaves, and its attachment measure
is invariant under exchanging root and summed leaf.
Fresh-block averaging is exact as well. Since \(a\)
is unique in its observed block, symmetry gives
\[
 \sum_c\gamma_{ac}\nu[(Y_c-Y_a)P_a]
   =\frac12\sum_c\gamma_{ac}
                            \nu[(Y_c-Y_a)(P_a-P_c)].
\]
The diagonal vanishes. At a distinct split of scale \(v\),
the allocation has total variation \(O(v)\) per dyad,
and \(\gamma_{ac}=O(v)\). The projected difference
costs a saved \(x^3/(Wv)\).

For the test difference, fix its deterministic replica
tree and write the factors not involving \(a,c\) as \(H\).
Let \(o\) be the exterior endpoint of an incoming edge,
whose split from the target trunk is \(j\in\{i,*\}\),
and put \(s_j=x\) or \(W\), respectively.
Set
\[
                         \mathcal F=\mathcal F_*^a,
\]
where \(\mathcal F_*^a\) is the sigma-field generated
by the complete planted prefix on the target trunk
through \(*\), including its common initial field.
No exterior continuation after its departure from
this trunk, no exterior terminal spin, and no target
innovation after \(*\) is included in \(\mathcal F\).
In particular \(X_*\) is \(\mathcal F\)-measurable.

All leaves used by \(H\) lie outside the current target
block and leave its trunk by \(*\). The product of
normalized planted transition kernels therefore makes
their joint continuation, together with \(y_o\),
independent of the joint future of \(a,c\), conditional
on \(\mathcal F\). Average this entire exterior joint
law and define
\[
 V=\mathbf E[Hy_o\mid\mathcal F],
       \qquad A=\mathbf E[y_oy_o^{\mathsf T}\mid\mathcal F].
\]
The projection powers at \(a,c\) are kept outside this
average. Explicitly, for any integrable scalar function
\(Z\) of their continuations and of \(\mathcal F\),
\[
 \nu[ZH\,(y_a-y_c)\cdot y_o]
             =\nu[Z\,(y_a-y_c)\cdot V].
\]
This identity applies with \(Z\) containing
\(Y_a-Y_c\) and the projection powers produced by
telescoping \(P_a-P_c\).

If \(j=*\), the unweighted marginal of \(y_o\)
conditional on \(\mathcal F\) is its ordinary
continuation from that prefix. If \(j=i<*\),
the extra target-trunk revelation between \(i\)
and \(*\) is independent of the exterior continuation
given \(\mathcal F_i^a\), so the same marginal
is the ordinary continuation from \(i\).
Consequently, for two independent copies \(o,o'\)
of this conditional marginal,
\[
 \|A\|_{\rm op}\le\|A\|_{\rm HS}
       =\bigl(\mathbf E[R_{oo'}^2\mid\mathcal F]\bigr)^{1/2},
\]
and their unconditional overlap has the split-\(j\)
marginal. Directional conditional Cauchy--Schwarz gives
\[
 \|V\|^2\le\mathbf E[H^2\mid\mathcal F]\,\|A\|_{\rm op}.
\]
Thus conditional Jensen and H\"older imply, for fixed
\(p\ge2\),
\[
 \|V\|_p\le\|H\|_{2p}\|R_j\|_p^{1/2}
                       \le C_p p_H\sqrt{s_j},
\]
where \(p_H\) is the raw scale of \(H\).
Correlations among \(H\) and the exterior spin are
permitted; none of their terminal observations
have been conditioned upon.

The vector \(V\) is measurable before the regular
innovation interval, which is chosen after \(*\)
and before the high split of \(a,c\).
Their common projection at the interval's end
cancels in \(V\cdot(y_a-y_c)\), so \eqref{eq:13}
applies to its two residuals with this fixed anchor.
Apply that estimate before H\"older with the remaining
target projection powers; those powers are neither
part of the anchor nor conditioned upon.
Terms where a projection power changes instead
use the raw incoming-edge bound and telescope
the powers of \(Y\). Together these operations
bound the test difference, in units of \(p_0\),
by a saved
\[
 \frac1v\left(\frac{x^3}{W^2}
                          +\frac{x^{5/2}}{W^{3/2}}\right).
\]
The first term covers projection powers and star incoming
edges; the second covers an incoming \(i\) edge.
Multiplying by the projected difference and the two
factors \(v\) gives, per dyad,
\[
 n^{-\kappa}p_0
  \left(\frac{x^6}{W^3}+\frac{x^{11/2}}{W^{5/2}}\right)
                           \le Cn^{-\kappa}p_0Wx^2.
\]
The saving absorbs the logarithmic number of dyads.

For the \(G_{ac}\) replacement one projection estimate
suffices. Allocate \(c\) first; its \(O(v)\) weight
cancels the denominator in \(x^3/(Wv)\).
The high centered error in \(G\) adds a saved \(x^2/W\).
If its split equals that of \(a,c\), use \eqref{eq:18}
over this same dyad; otherwise allocate that error pair
next and use Lemma~\ref{lem:found-split-density}. Reordering these
allocations retains the full scalar topology coefficient, by the
labeled product rule in Proposition~\ref{prop:found-replica}; that
coefficient is bounded here and is not collapsed as a free row.
The result is
a saved \(p_0x^5/W^2\le p_0Wx^2\).
All bounds are uniform in the target band.
\end{proof}

Apply \eqref{eq:24} with \(P^\circ=P-\nu[P]\).
Constants against \(M\) alone vanish, so the replacements give
\begin{equation}\label{eq:25}\tag{25}
 \begin{split}
 (1-b_*\chi^2)\nu[Y_aP^\circ]
   ={}&\chi K_*\{\mathcal L(P,a)+\mathcal L(P,b)\}\\
     &+b_*\left\{\chi\mathcal L(P^\circ Y_a,a)
       +\sum_c\gamma_{ac}\mathcal L(P^\circ\Delta_{bc},b)\right\}
       +\mathcal E(P;a,b).
 \end{split}
\end{equation}
Here the row identity and projected replacements give
\[
 |\mathcal E(P;a,b)|\le
 C_NWx^2p_0\{\eta_{N,n}+(D_i/x)^{\vartheta_N}\},
\]
with \(\eta_{N,n}\to0\), \(\vartheta_N>0\), and no \(D_i\)
term for star-only tests.
After switching to \(b\), its single incoming edge is from
\(c\) in the former target block, as required by
Lemma~\ref{lem:cavity-family}.

\begin{lemma}[Finite-depth longitudinal cancellation]\label{lem:cavity-recursion}
Put \(J_0=T\chi^2-1\). Every admissible fixed-degree test satisfies,
uniformly over initial targets \(i\) in the band of
Proposition~\ref{prop:cavity-closure},
\[
 |\mathcal L(P,a)|\le
 Cx^2p_0\{\eta_n+(D_i/x)^\vartheta+\eta_n|J_0|/W^2\},
                       \qquad \eta_n\to0,\quad\vartheta>0.
\]
The \(D_i\) term is absent for star-only tests.
Consequently, for such tests and along moving initial targets with
\(D_i=o(x)\),
\begin{equation}\label{eq:26}\tag{26}
       \mathcal L(P,a)=o\bigl(x^2p_0(1+|J_0|/W^2)\bigr).
\end{equation}
\end{lemma}

\begin{proof}
Equations \eqref{eq:4} and \eqref{eq:6} give
\[
 b_*=T+O(W^2),\qquad K_*=TS_*+O(W^3),\qquad \chi K_*\asymp W.
\]
The good-split bound and raw higher moments yield
\[
 \frac{|(1-b_*\chi^2)\nu[Y_aP^\circ]|}{\chi K_*}
            \le\eta_nx^2p_0(1+|J_0|/W^2).
\]
Any fixed interpolation loss in the polynomial \(D_*\)
saving is included in \(\eta_n\).
To display one generation of the recursion, put
\[
 \mathcal C(P;a,b)=
 \frac{b_*}{K_*}\mathcal L(P^\circ Y_a,a)
 +\frac{b_*}{\chi K_*}
        \sum_c\gamma_{ac}\mathcal L(P^\circ\Delta_{bc},b).
\]
After division by \(\chi K_*\), Equation~\eqref{eq:25} reads
\[
 \mathcal L(P,a)+\mathcal L(P,b)
       =\varepsilon(P;a,b)-\mathcal C(P;a,b),
\]
where, uniformly over the allowed target band,
\[
 |\varepsilon(P;a,b)|\le C_Nx^2p_0
   \{\eta_n+(D_i/x)^\vartheta+\eta_n|J_0|/W^2\}.
\]
The target term is absent for star-only tests. Every term in
\(\mathcal C\) adds one star factor, with coefficient \(O(1/W)\).

Now choose two fresh roots \(u,v\) at \(*\), neither in \(P\).
Pruning the unused fresh root preserves the normalized planted law;
relabeling therefore identifies each of
\(\mathcal L(P,u)\) and \(\mathcal L(P,v)\) with
\(\mathcal L(P,b)\). Apply the same identity to \(u,v\), and
subtract half of it from the identity for \(a,b\). This gives
\[
 \mathcal L(P,a)
 =\varepsilon(P;a,b)-\tfrac12\varepsilon(P;u,v)
  -\mathcal C(P;a,b)+\tfrac12\mathcal C(P;u,v).
\]
Every insertion on the right now has one more star factor than
the insertion on the left. This is the recursion we iterate.

Fix a depth \(k\), independent of \(n\).
At depth \(j\) the natural scale is \(p_0W^j\),
and accumulated coefficients are at most
\(C_jW^{-j}\). Lemma~\ref{lem:cavity-family}
ensures admissibility after every switch, even
with old high subtrees. Ancestor coefficients
are fixed during each new interpolation.
Row weights have total variation bounded
in terms of \(j,N\). Thus intermediate errors
are bounded by
\[
 C_{k,N}x^2p_0\{\eta_{k,N,n}+(D_i/x)^{\vartheta_{k,N}}
                                  +\eta_{k,N,n}|J_0|/W^2\}.
\]
There are finitely many such errors.

At a terminal node there are exactly \(k\)
new marked star factors, centered overlaps or
projections. H\"older at its required fixed
order and marginalization identify each with
a star pair; \eqref{eq:17} bounds each by
\(C_kn^{-c_2}W\), with \(c_2\) independent
of that order. Initial factors use raw scales.
The high insertion has bounded fixed moments
and finite-label total variation. Including
the \(W^{-k}\) normalization, the terminal
bound is \(C_{k,N}p_0n^{-kc_2}\).
Centering can send marked factors into
earlier separate expectations; applying this
argument to each expectation preserves their
total number \(k\). Choose \(k\) fixed with
\(kc_2>1/3\). Since \(x\ge n^{-1/6+m}\),
the terminal bound is \(o(x^2p_0)\).
Only then fix the degrees and H\"older orders.
Summing the finite recursion proves the lemma.
\end{proof}

\begin{lemma}[Susceptibility calibration]\label{lem:cavity-calibration}
One has \(J_0=O(W^2)\). Therefore
\begin{equation}\label{eq:27}\tag{27}
                         \mathcal L(P,a)=o(x^2p_0)
\end{equation}
for admissible tests along moving targets with \(D_i=o(x)\),
and uniformly
\[
 |\mathcal L(P,a)|\le Cx^2p_0\{\eta_n+(D_i/x)^\vartheta\}.
\]
\end{lemma}

\begin{proof}
Use \eqref{eq:24} with \(P=1\). The linear centered term
and the constant \(K_*\) parts of the high corrections
vanish by marginalization. By
Lemma~\ref{lem:cavity-replacements}, remaining terms are
insertions against one star factor. Equation \eqref{eq:26}
with \(p_0=W\) bounds them by
\(o(Wx^2(1+|J_0|/W^2))\).

The scalar low expansion is
\(\Gamma(K_*)=\chi^2K_*+O(W^3)\).
Its second coefficient is a two-edge path
through a third block. With no other observations
the block is unobserved, gaining \(O(W)\),
and its two deterministic low factors are
\(O(W^2)\). The third derivative is
\(O(W^3)\) by bounded finite-label variation.
Consequently
\[
 S_*=\chi^2TS_*+O(W^3)+o(Wx^2(1+|J_0|/W^2)).
\]
Divide by \(S_*\asymp W\) to obtain
\[
 |J_0|\le CW^2+\eta_nx^2+\eta_n|J_0|x^2/W^2.
\]
Since \(x\le W\), absorb the last term.
Lemma~\ref{lem:cavity-recursion} now gives
\eqref{eq:27} and its uniform version.
\end{proof}

\subsection{Return to the target cut}

\begin{proof}[Proof of Proposition~\ref{prop:cavity-closure}]
Choose a regular continuity cut \(U_x\asymp C_2x\)
above the entire fixed target band. Repeat the
mean expansion of Lemma~\ref{lem:cavity-row-identity}
with split \(i\) and \(P=1\), also allowing \(W\)
bounded below. The scalar-law and block conventions in this repetition
use \(U_x\) in place of \(U\). All low factors cost \(O_p(x)\).
Above \(U_x\), the dyadic bounds in \eqref{eq:18} are
polynomially smaller than \(x\). Indeed, on the intermediate
dyads and the high dyads, respectively,
\[
 \frac{x^2}{(W^2v)^{1/3}}\lesssim x
       \quad(v\ge C_2x,\ W\ge x),
 \qquad \frac{x^2}{W}\le x.
\]
These are bounds after split integration; no pointwise high-split
estimate is needed. A third-remainder atom at \(i\)
costs \(O(x^2D_i)\); an integrated low error
costs \(x\int_0^{C_2x}D=o(x^3)\).
High Taylor remainders cost two saved high
errors times \(x\), or one saved high error
times \(x^2\), and hence are \(o(x^3)\).
The centered first-coefficient
low restoration has zero first derivative
and second cost \(O(x^2D_i)\).
At second low order every third block is
unobserved; its departure density or \(O(x)\)
entry coefficient gives the same product
difference bounds. This exhausts the earlier
ledger at scale \(x\), with remainder
\[
                       Cx^3\{\eta_n+(D_i/x)^\vartheta\}.
\]

The linear centered term and deterministic high
means vanish. In the remaining \(G_{ac}\) term
exchange \(a,c\) in the sum and average the
opposite fresh branch. Let \(M_a^{(x)},\chi_x\)
denote the insertion and scalar row at \(U_x\),
and \(Y_a^{(i)}=X_i\cdot y_a-S_i\).
The partner term is identical. The surviving
form is therefore
\[
 S_i-\Gamma(K_i)=
       2b_i\chi_x\,\nu[Y_a^{(i)}M_a^{(x)}]
                                      +\text{remainder}.
\]
No test depending on \(a\) preceded this mean
expansion, so this exchange is exact.

If \(W\) is bounded below along a subsequence,
\eqref{eq:18}--\eqref{eq:19} give
\(\int_{U_x}^1D_u\,du=o(x^2)\):
up to \(C_3x\) use \eqref{eq:19}, from
\(C_3x\) to a fixed multiple of \(W\) use
the intermediate dyadic bound of
\eqref{eq:18}, and above use its high-band
bound. Take \(C_3\) sufficiently large and
fixed. Allocate the error pair first,
relative only to \(a\) in its block, giving
a bounded density. The insertion costs
\(O(x)\int D=o(x^3)\).

If \(W=o(1)\), choose the calibration cut
\(U\) above \(U_x\), adjusting its fixed
multiple also when \(W\asymp x\).
Remove error edges splitting between the
cuts at cost
\[
                        Cx\int_{U_x}^{U}D_u\,du=o(x^3).
\]
For remaining edges replace the scalar
connected kernel at \(U_x\) by the one
with revelation off up to \(U\).
Each bounded-spin scalar moment changes
by \(O(W)\): its covariance change below
\(U\) is \(O(W)\), and \eqref{eq:2}
has bounded finite-label variation.
The insertion error is at most
\[
                             CxW\int_U^1D_u\,du=o(x^3)
\]
by the saved \(x^2/W\) high bound.
No arbitrary high-atom estimate is needed:
allocate the error pair first relative
to the sole old root \(a\) in this block,
and then allocate the summed leaf \(c\).

For the new kernel, independent block
flips force \(c\) into \(a\)'s calibration
block. The error endpoints share a block
because their split is above \(U\).
If it is different from the block of
\(a,c\), the connected scalar covariance
vanishes. Thus the remaining operator is
exactly \(M_a\) of \eqref{eq:23}.
The test \(Y_a^{(i)}\) belongs to our family
by its fresh-tooth representation and
has one initial factor of scale \(x\).
Equation \eqref{eq:27} proves
\(\nu[Y_a^{(i)}M_a]=o(x^3)\), with its
uniform target modulus.

Cuts and integrated estimates are chosen
once for the full fixed target band.
After fixing the recursion depth only
finitely many H\"older orders are used.
Target dependence is a finite collection
of positive powers of \(D_i/x\); their
minimum is a common \(\vartheta>0\).
All remaining errors give a common
\(\eta_n\to0\).

For omitted lower cuts suppose \(W=O(x)\), hence
\(W\asymp x\), and choose all regular cuts at the
required fixed multiples of \(x\), uniformly in
the omitted target. First perform the star-only
calibration: it gives
\[
                 J_0=O(x^2),\qquad \chi K_*\asymp x,
\]
using \eqref{eq:4} only in the regular calibration
region. At every omitted \(i\), the global raw
bounds and conditional Jensen give
\[
 \|R_i\|_p+|S_i|+D_i+\|Y_a^{(i)}\|_p\le C_px,
                  \qquad \|X_i\|_p\le C_p\sqrt x .
\]
Monotonicity from the regular region also gives
\(h_i,K_i\le Cx\).

Enlarge the admissible tests by allowing their
one initial edge to split at any omitted \(i\),
retaining its natural scale \(x\). The row-identity
proof remains uniform for this class: an old-\(i\)
error costs at most \(C_Nx^2(D_i+1/n)p_0\);
an integrated low error costs
\(C_Nx p_0\int_0^{Cx}(D_u+1/n)\,du\).
The observed-third-block baseline at \(i\) has
product difference \(O_p(x^2)\) and gains a
cut mass \(O(x)\). The other remainder classes
retain their previous bounds. Thus the remainder
in \eqref{eq:24} is \(O_N(x^3p_0)\).
All projection and resampling intervals lie in
the regular region at masses at least \(Cx\).
For an initial incoming edge the weighted-anchor
bound uses only its split-\(i\) raw moments and
is \(O_p(p_H\sqrt x)\). Its projection denominator
is the regular innovation mass, never \(u_i\).
Finite-label attachment variation is bounded
independently of \(u_i\); zero splits follow
by the same degenerate-tree conventions and
continuity.

Use \eqref{eq:25} with this enlarged class,
after the star-only calibration, and divide
by \(\chi K_*\asymp x\). Its left side and
remainder are \(O_N(x^2p_0)\).
The two-fresh-root equation removes the
same-degree nuisance as before. At depth \(j\),
the test scale is \(p_0x^j\) and the accumulated
coefficient is at most \(C_jx^{-j}\), so
each intermediate error is \(O_j(x^2p_0)\).
The \(k\) marked star factors bound terminal
terms by \(C_kp_0n^{-kc_2}\), also when centering
places some marks in separate expectations.
Choose fixed \(kc_2>1/3\); then
\(n^{-kc_2}=o(x^2)\). The finite recursion proves
\[
                         |\mathcal L(P,a)|\le C_Nx^2p_0.
\]

In particular \(P=Y_a^{(i)}\) has scale \(x\),
so \(|\nu[Y_a^{(i)}M_a]|\le Cx^3\).
The preceding cutoff and scalar-kernel changes
still cost \(o(x^3)\), uniformly in \(i\).
The target mean identity therefore gives
\[
 S_i-\Gamma(K_i)
      =2b_i\chi_x\,\nu[Y_a^{(i)}M_a]+O(x^3)=O(x^3),
\]
since \(b_i,\chi_x\) are bounded. Here
\(K_i=h_i+b_iS_i\) is always the actual changed
clock: it is never replaced by \(TS_i\), so
no term \((b_i-T)S_i\) occurs. No endpoint
hypothesis is used at the omitted cut.
Finally, if the omitted mass is \(u_0=o(x)\),
its contribution to the integrated \(D\)
bound is at most \(Cx u_0=o(x^2)\).

Finally let \(d\mu_n=du/x\) on a mass band
of \eqref{eq:19}, or \(d\mu_n=dr/x\) on
its radius band, restricted in each case
to the indicated set. These measures have
uniformly bounded total mass. Equation
\eqref{eq:19} and the raw bound imply
\[
 \int(D_i/x)^2\,d\mu_n(i)\longrightarrow0,\qquad D_i/x\le C.
\]
Decrease \(\vartheta\) to at most 2.
Jensen gives
\[
 \int\frac{|S_i-\Gamma(K_i)|}{x^3}\,d\mu_n(i)
 \le C\eta_n+
 C\left(\int(D_i/x)^2\,d\mu_n(i)\right)^{\vartheta/2}
 \longrightarrow0.
\]
This proves both normalized \(L^1\) claims.
\end{proof}
